% Dueño: B. A21-A23, A29, A33: tasa de recuperación y archivos sin codificar.
\subsection{Tasa de recuperación y rechazo de señales sin codificar}
\label{sec:recuperacion}

Para caracterizar la tasa de recuperación en función de la atenuación del eco y del género musical, se evaluaron los cuatro fragmentos de prueba, codificados con bits al azar mediante el prototipo en Python. Con la configuración por defecto ($a = \SI{30}{\percent}$), la tasa media global alcanza \SI{98.31}{\percent} (\SI{97.57}{\percent} en la clásica y en el \textit{ragtime}, \SI{98.13}{\percent} en el \textit{jazz} y \SI{100.00}{\percent} en electrónica), superando el \SI{90}{\percent}. Con \SI{10}{\percent} de atenuación las tres grabaciones reales quedan cerca del \SI{80}{\percent}, mientras que la electrónica sintética, de espectro estable, se mantiene sobre el \SI{99}{\percent}.


Con $a = \SI{30}{\percent}$ y otros bits al azar, una ventana mayor reduce la tasa binaria, de \SI{66.82}{bit\per\second} con $K=10$ a \SI{26.73}{bit\per\second} con $K=25$ para $(750, 1500)\,\si{\micro\second}$, pero la recuperación no crece de forma monótona con $K$: para ese par va de \SI{95.53}{\percent} con $K=10$ a \SI{91.69}{\percent} con $K=15$ y \SI{97.94}{\percent} con $K=25$. 


Finalmente, ante los cuatro archivos originales sin codificar, el filtro de actividad global produce exactamente cero bits emitidos y cero metadatos falsos.

